% Storyline 1, paragraph 5 (round 2): findings; one sentence carries the Storyline 2 perspective.
We find that the allocation of inertia between the two directions, rather than its total amount, determines whether private observers agree with one another and still distinguish cooperators from defectors. Because observers share no information, this agreement emerges from how each individual revises its own judgements. The beneficial direction depends on the norm. Compact architectures, selected solely by these reputational criteria, make cooperation stable and easy to reach under each of four classical norms (Scoring, Shunning, Stern Judging and Simple Standing). At the same depth, inertia placed on the wrong path can leave no cooperation at all. We further show that the tolerance that protects justified punishers can also shelter opportunists who defect only occasionally, a vulnerability that the aggregation and score-based rules we test share. Allowing warning states to carry partial consequences closes this loophole for the opportunistic strategies we consider, while cooperation remains high. Reputation-based cooperation therefore depends not only on which behaviours a society judges as good, but also on how readily it revises those judgements in each direction and on how it treats individuals before a judgement changes.
